\documentclass[10.5pt,a4paper]{article}
\usepackage[T1]{fontenc}
\usepackage[utf8]{inputenc}
\usepackage[english,italian]{babel}
\usepackage{lmodern,microtype}
\usepackage[a4paper,margin=1.85cm,headheight=15pt]{geometry}
\usepackage{enumitem,tabularx,array,longtable,booktabs}
\usepackage{xcolor,hyperref,xurl,fancyhdr,titlesec,framed}
\definecolor{ddtadark}{RGB}{28,38,52}
\definecolor{ddtablue}{RGB}{42,88,135}
\definecolor{ddtagray}{RGB}{244,246,248}
\definecolor{ddtagreen}{RGB}{48,111,78}
\definecolor{ddtaorange}{RGB}{148,91,25}
\hypersetup{colorlinks=true,linkcolor=ddtablue,urlcolor=ddtablue,
  pdftitle={SRC-0034 - scheda compilata BA0-R0},
  pdfauthor={Documentation-Driven Threat Analysis research workspace}}
\pagestyle{fancy}\fancyhf{}\lhead{DDTA - SRC-0034}\rhead{BA0-R0 bridge rereading}\cfoot{\thepage}
\titleformat{\section}{\large\bfseries\color{ddtadark}}{}{0pt}{}
\setlist[itemize]{leftmargin=1.4em,itemsep=0.18em,topsep=0.25em}
\setlength{\parindent}{0pt}\setlength{\parskip}{0.33em}\setlength{\emergencystretch}{2em}
\newenvironment{factbox}[1]{\def\FrameCommand{\fcolorbox{ddtablue}{ddtagray}}\MakeFramed{\advance\hsize-2\width\FrameRestore}\noindent\textbf{#1}\par\smallskip}{\endMakeFramed}
\newenvironment{gatebox}[1]{\def\FrameCommand{\fcolorbox{ddtagreen}{ddtagray}}\MakeFramed{\advance\hsize-2\width\FrameRestore}\noindent\textbf{#1}\par\smallskip}{\endMakeFramed}
\newenvironment{cautionbox}[1]{\def\FrameCommand{\fcolorbox{ddtaorange}{ddtagray}}\MakeFramed{\advance\hsize-2\width\FrameRestore}\noindent\textbf{#1}\par\smallskip}{\endMakeFramed}
\begin{document}

\begin{center}
{\LARGE\bfseries DDTA - Scheda di lettura compilata}\par
\vspace{0.2em}
{\large SRC-0034 - BA0-R0 bridge rereading}\par
\vspace{0.3em}
{\small Rilettura mirata prima del systems-modeling prior-art corpus. Non sostituisce la source note o l'excerpt ledger esistenti.}
\end{center}

\begin{factbox}{Identita bibliografica e accesso}
\begin{tabularx}{\linewidth}{>{\bfseries}p{3.4cm}X}
Source ID & SRC-0034\\
Citation key & \texttt{souza2019architecturalmodels}\\
Titolo & \textit{Deriving architectural models from requirements specifications: A systematic mapping study}\\
Autori & Eric Souza; Ana Moreira; Miguel Goul\~ao\\
Anno / venue & 2019; Information and Software Technology 109, 26--39\\
Tipo di fonte & secondary review / systematic mapping study\\
Stable identifier & DOI \texttt{10.1016/j.infsof.2019.01.004}\\
Access version & publisher-formatted institutional repository copy\\
Local file & \texttt{SRC-0034-souza-moreira-goulao-2019.pdf}\\
Retrieved & 2026-08-06\\
Local SHA-256 & \texttt{7560a403216ba0df46be238bf7bf075ef016be2c028182aa6d8a4ee28fd22adf}\\
Verification state & verified\\
\end{tabularx}
\end{factbox}

\section{1. Problema di ricerca e contributo}
La fonte mappa i metodi proposti per derivare modelli di software architecture da requirements specifications. Il passaggio osservato non e una semplice traduzione sintattica: separa comprensione del problema, ricerca di una soluzione ed evaluation della soluzione, con alternative e trade-off architetturali.

Il corpus finale comprende 39 primary studies selezionati da 2575 candidati. La review organizza i metodi secondo quattro dimensioni NIMSAD: context, benefits to users, method content e validation. Il valore centrale per BA0-R e il confine requirements-to-architecture, non l'identificazione di un metamodel minimo universale.

\section{2. Artefatti iniziali e rappresentazione del sistema}
Tutti i 39 studi trattano functional requirements; 30 includono anche NFR. Gli input sono eterogenei: textual requirements (20), goal-oriented models (8), unspecified models (6), feature models (2), problem frames (2) e sequence diagram (1). Anche nel sottogruppo testuale compaiono use cases, intermediate models e clustering.

La classificazione DDTA gia registrata e \texttt{requirements\_primary\_with\_model\_mediation}: i requirements sono la sorgente principale, ma la derivazione passa spesso attraverso rappresentazioni strutturate o decisioni dell'architetto.

Le rappresentazioni di output sono a loro volta eterogenee. Diciannove studi non nominano esplicitamente le views; gli altri usano 28 nomi di view per 80 istanze. Lo stesso nome puo indicare semantiche differenti e nomi diversi possono indicare scopi simili. Tredici studi usano ADL standard (inclusi UML, SysML o AADL), 13 notazioni non standard e 13 non identificano la notazione.

\begin{cautionbox}{BA0-R0 consequence}
La fonte giustifica una pressione contro \texttt{Actor + Asset + Interaction} come core sufficiente, ma non autorizza ancora \texttt{Channel}, \texttt{Store}, \texttt{Interface}, \texttt{Boundary} o altri concetti come BAE first-class. Il corpus e troppo eterogeneo per derivare da solo un minimo comune.
\end{cautionbox}

\section{3. Traceability, change, automation e ruolo umano}
Il pattern di automazione osservato e representation-specific: una specifica rappresentazione dei requirements viene trasformata in uno specifico candidate architectural model, seguito da manual analysis and improvement. La review riporta 18 studi parzialmente automatizzati, 8 non automatizzati, 13 senza tool menzionato e nessun approccio end-to-end completamente automatizzato nel corpus storico.

Le decisioni architetturali dipendono comunque da esperienza e intuizione dell'architetto. I compiti umani comprendono identificazione degli architecturally significant requirements, interpretazione di NFR/domain constraints, generazione e confronto di alternative, scelta di pattern/styles/tactics, trade-off, correzione degli artefatti generati ed evaluation della satisfaction.

La traceability requirements-to-architecture e dichiarata come beneficio in un solo studio. Longitudinal synchronization, stale detection e change-impact propagation non sono valutati. La fonte richiama inoltre il rischio di architectural knowledge vaporization: il rationale puo perdersi anche quando la struttura corrente resta visibile.

\section{4. Findings, limiti e relazione con DDTA}
La review segnala dipendenza da tacit knowledge, terminologia non standard, debole supporto ai novizi, automazione parziale e validation insufficiente. Ventuno studi non dichiarano come valutare l'architettura risultante rispetto ai requirements; quindi:

\begin{center}\ttfamily
architecture generated $\neq$ requirements satisfied\\
$\neq$ alternatives compared $\neq$ trade-offs justified $\neq$ architecture validated
\end{center}

Per DDTA la fonte sostiene il bisogno di separare source evidence, candidate representation, decision/review e accepted representation. Rafforza inoltre l'idea che una view sia una proiezione con semantica dichiarata, non necessariamente il modello canonico.

\begin{gatebox}{Novelty challenge BA0-R0}
SRC-0034 mostra sovrapposizione con DDTA nel passaggio requirements-to-model/architecture, ma non una pipeline end-to-end equivalente a governed documentation $\rightarrow$ methodology-neutral analyzable representation $\rightarrow$ replaceable security analysis $\rightarrow$ reviewed findings $\rightarrow$ governed change/re-analysis. Il reopen trigger di Chapter 3 non e soddisfatto da questa fonte.
\end{gatebox}

\section{Metodo di ricerca e corpus}
Ricerca su IEEE Xplore, ACM Digital Library, ScienceDirect e SpringerLink, con forward snowballing ed expert suggestions. Il protocollo e stato rivisto da tre external reviewers, piloted e controllato con Fabbri's checklist. Il quality gate combina allineamento al problema, formalizzazione del metodo, architecture description, evaluation, limitations, venue ranking e citation information.

E presente una reporting inconsistency conservata nella note DDTA: Table 6 non si riconcilia aritmeticamente nella colonna ScienceDirect; il totale finale di 39 usa il valore mostrato dagli autori senza correzione arbitraria.

\section{Validita temporale}
\begin{tabularx}{\linewidth}{>{\bfseries}p{4.1cm}X}
Contribution type & historical systematic mapping study\\
Primary-study horizon & studi disponibili fino circa al 2016; data di search non esplicitata nella note\\
Time sensitivity & medium per gap concettuali; high per tool/automation maturity\\
Ancora applicabile & necessita di decisioni esplicite, multi-view semantics, traceability ed evaluation\\
Potenzialmente obsoleto & livello di automazione, supporto tool, industrial adoption, assenza di LLM moderni\\
Corroborazione richiesta & necessaria per qualsiasi claim sullo stato 2026 di tool/LLM\\
Uso consentito & historical boundary, taxonomy/governance implications\\
Overclaim vietato & current inventory o impossibilita di automazione piu recente\\
\end{tabularx}

\section{Relazione con DDTA e implicazioni BA0-R0}
\textbf{Direttamente supportato}
\begin{itemize}
\item Requirements testuali o modellati possono iniziare la derivazione architetturale.
\item La derivazione include analisi, sintesi e evaluation e non determina una sola soluzione.
\item Multiple views e ADL sono usate con terminologia eterogenea.
\item L'automazione parziale non elimina judgment e review architetturale.
\item Traceability e requirements-satisfaction validation sono deboli nel corpus mappato.
\end{itemize}

\textbf{Non supportato}
\begin{itemize}
\item Un set minimo di primitive BAE per DDTA.
\item Una canonical neutral representation comune a tutti gli approcci.
\item Current 2026 LLM/tool maturity.
\item Derivazione deterministica o correctness automatica dell'architettura.
\end{itemize}

\textbf{Pressioni da portare a NASA/SysML/KerML}
\begin{itemize}
\item distinguere element, relation, property e view;
\item verificare se behavior, information/flow, interface, channel/link, persistence, state/mode e boundary sono responsabilita semantiche ricorrenti;
\item distinguere entity, role, relation e property;
\item verificare canonical model vs projection/view;
\item mantenere explicit review authority tra extraction e accepted representation.
\end{itemize}

\section{Citation-ready locations usate nella rilettura}
\begin{longtable}{p{1.8cm}p{1.15cm}p{1.3cm}p{3.2cm}p{7.1cm}}
\toprule
Excerpt & Tipo & PDF p. & Sezione & Proposizione supportata\\
\midrule
\endhead
EX-0002 & P & 1 & Introduction & Architecture derivation distingue problem understanding, solution finding ed evaluation.\\
EX-0012 & P & 5 & 3.2 Context & Gli input includono requirements testuali e diversi modelli strutturati.\\
EX-0017 & P & 6 & 3.4 Content & Le architectural views hanno naming e semantiche eterogenee.\\
EX-0019 & P & 8 & 3.4 Content & UML/SysML/AADL sono solo parte delle notazioni usate.\\
EX-0020 & P & 8 & 3.4 Content & Nessun fully automated end-to-end approach nel corpus; serve manual improvement.\\
EX-0023 & P & 10 & 3.5 Validation & Piu di meta dei metodi non esplicita evaluation contro i requirements.\\
EX-0027 & P & 12 & 7 Related work & Architectural knowledge vaporization durante l'evoluzione.\\
EX-0028 & P & 12 & Conclusions & Derivazione ancora expert-led e incompletamente validata.\\
\bottomrule
\end{longtable}

\section{Follow-up sources per BA0-R}
Questa rilettura non registra nuove fonti. Conferma soltanto la necessita del seed corpus gia pianificato: NASA-HDBK-1009A, SysML 2.0, KerML 1.0, NASA Systems Engineering Handbook, ISO/IEC/IEEE 42010, AADL, ArchiMate e survey MBSE. Le bibliografie di tali fonti saranno minate separatamente e ogni candidato dovra superare verifica bibliografica/access gate prima di ricevere un nuovo SOURCE-ID.

\begin{gatebox}{BA0-R0 status}
\textbf{CLOSED.} Exact BAE primitives remain OPEN. BA0 remains PAUSED and BA1 NOT STARTED. Next research microstep: BA0-R1, NASA-HDBK-1009A.
\end{gatebox}

\vfill
{\footnotesize\textbf{Evidence governance.} Questa scheda riassume esclusivamente contenuti gia registrati in \texttt{literature/notes/SRC-0034.md} e \texttt{literature/excerpts/SRC-0034.excerpts.yml}. Per citazioni e claim di tesi, usare sempre quei record e le posizioni esatte nella fonte.}
\end{document}
